\documentclass[10.5pt]{article}
\usepackage[utf8]{inputenc}
\usepackage[margin=0.8in,top=0.6in,bottom=0.65in]{geometry}
\usepackage{titlesec}
\usepackage{parskip}
\usepackage{fancyhdr}
\usepackage{amsmath, amssymb}
\usepackage[round,sort&compress]{natbib}
\usepackage{hyperref}
\renewcommand{\bibsection}{\section*{References}}
\setlength{\bibsep}{0pt}

\titlespacing*{\section}{0pt}{5pt}{2pt}
\setlength{\parindent}{0pt}
\setlength{\parskip}{2.5pt}

\pagestyle{fancy}
\fancyhf{}
\lhead{\footnotesize STAT8101}
\chead{\footnotesize \href{http://www.cs.columbia.edu/~blei/teaching.html}{Invariance, Causality, and Applications}}
\rhead{\footnotesize \href{http://www.cs.columbia.edu/~blei/}{David M. Blei} (\texttt{db2128})}
\cfoot{\thepage}

\title{\vspace{-2.5em}\textbf{Reader Report: Week 06} \\ \vspace{0.1em} \large Probabilistic Invariance: VI-BIP and the Core Theorems}
\author{\href{https://github.com/dhardestylewis}{Daniel Hardesty Lewis} (\texttt{dl3645})}
\date{\vspace{-0.5em}2026 Spring}

\begin{document}

\maketitle
\thispagestyle{fancy}
\vspace{-1.2em}

\section*{Summary}

\citet{wu2025bayesian} place a prior over binary selectors $z\in\{0,1\}^p$, where $z_j=1$ indexes a feature included in the invariant set $S^*$ \citep{peters2016icp}.
The posterior
$p(z\mid\mathcal{D})\propto p(z)\prod_{e,i}\hat{g}(Y^e_i\mid X^z_{ei})/\hat{p}_e(Y^e_i\mid X^z_{ei})$
rewards selectors for which the pooled conditional $\hat{g}$ dominates each local conditional $\hat{p}_e$, where $\hat{p}_e$ is fit separately per environment $e$.
Consistency (Thm~1) requires $\mu_{\min}>0$, defined as the minimum KL divergence between $\hat{g}$ and any $\hat{p}_e$ over non-invariant features \citep[Def.~1]{wu2025bayesian}.
Contraction (Thm~2) gives posterior mass on $\{z\neq z^*\}$ scaling as $O(R\,e^{-\kappa nE\mu_{\min}})$, where $R=2^p$ counts hypotheses, $E$ is the number of environments, and $\kappa$ is a universal constant.
VI-BIP \citep[Sec.~4]{wu2025bayesian} approximates the posterior via mean-field variational inference \citep{blei2017variational}, optimizing the evidence lower bound (ELBO) using the U2G gradient estimator to handle discrete $z$.

\section*{Critique: One Invariance Condition, Three Algorithms}

The most striking observation across this body of work is that \citet{peters2016icp}, \citet{fan2024negdro}, and \citet{wu2025bayesian} all enforce the invariance condition $P_e(Y\mid X_{S^*}) = \text{const}$ for all $e\in\mathcal{E}$, yet operationalize it through algorithms that produce different outputs: a confidence set, a point estimate, and a posterior distribution, respectively.

Peters et al.\ test invariance sequentially, building a confidence set $\hat{S}(\mathcal{E})$ that shrinks monotonically as environments accumulate \citep[Thm.~1]{peters2016icp}.
BIP's Thm~2 is the probabilistic counterpart: rearranging the contraction gives sample complexity $nE \gtrsim O(p/\mu_{\min})$, which is the per-environment requirement for the Peters test to have nontrivial power.
BIP thus does not replace ICP; it quantifies ICP's frequentist power curve in a Bayesian language.

NegDRO \citep{fan2024negdro} makes a third move.
Its minimax problem $\min_b\max_{w\in U(\gamma)}\sum_e w_e\hat{\mathbb{E}}[(Y^e{-}b^\top X^e)^2]$, where $U(\gamma)$ is the set of signed environment weights with total variation at most $\gamma$ \citep[Def.~1]{fan2024negdro}, is the distributionally robust framing of \citet{buhlmann2020invariance} pushed to an adversarial extreme.
BIP treats environments symmetrically in the product likelihood; NegDRO can amplify a single harsh environment pair via negative weights, making it more powerful when one intervention is known \emph{a priori} to be dominant.
The two methods have different assumptions about what the analyst knows before seeing the data, which is a structural incompatibility neither paper acknowledges.

A further issue concerns benchmarking. \citet[Sec.~5]{wu2025bayesian} validate VI-BIP on a single genomics dataset: 6{,}170 yeast genes across CRISPR knockout conditions where $\mu_{\min}>0$ holds by construction. The paper provides no evaluation in low-environment or observational settings, where $\mu_{\min}$ is unknown; whether the contraction advantage over ICP survives soft-intervention regimes is unresolved.

\section*{Questions}
\begin{enumerate}\setlength\itemsep{1pt}
  \item NegDRO's $\lambda$ and BIP's $\mu_{\min}$ both quantify environment heterogeneity but measure different objects: $\lambda$ is a spectral condition on second moments of $X$, while $\mu_{\min}$ is a KL condition on $Y|X$ conditionals. Is it possible to certify $\mu_{\min}>0$ from observational data without already knowing $z^*$, or does this require the hard interventions that make the genomics setting tractable?
  \item Wu et al.\ propose a size-restricted prior with parameter $p_{\max}$ to limit BIP's support to subsets of size at most $p_{\max}$. Could NegDRO's point estimate $\hat{\beta}^\gamma$ serve as the pre-screening filter, setting $z_j=0$ for features with $|\hat{\beta}^\gamma_j|<\varepsilon$ to reduce the effective model count from $2^p$ to $2^{p'}$, and does this avoid the double-dipping problem that arises from using the same data to form both the prior and the likelihood?
\end{enumerate}

\vspace{-0.2em}
{\scriptsize\setlength{\bibsep}{0pt}
\bibliographystyle{abbrvnat}
\bibliography{references}
}

\end{document}
